% GENERATED FILE. Edit canonical lessons, then run npm run generate:books.
% canonical-source-hash: fnv1a64:fea0c360b719b203
% canonical-lessons: ZH-C267-yanyuan, ZH-C267-jiemu, ZH-C267-mingxing, ZH-C267-xiaohua, ZH-W267-xiao-laugh, ZH-C267-haoxiao

\chapter{Actors and Jokes}
\label{ch:zh-actors-and-jokes}
\hlchaptermodality{267}

\begin{chapteropening}
\textbf{By the end of this chapter:} \emph{I can talk about actors, TV shows and stars, and jokes that are funny.}

Complete the last lesson of chapter 267: I can talk about actors, TV shows and stars, and jokes that are funny.
\end{chapteropening}

\section[yǎnyuán]{yǎnyuán — actor; performer}

\label{lesson:ZH-C267-yanyuan}



\begin{quote}
Before the new one: say the Mandarin for a girl, then the Mandarin for middle age.
\end{quote}

\subsection*{You'll want to know: yǎnyuán}
\emph{yǎnyuán} — \textquotedblleft{}actor; performer\textquotedblright{}.

\begin{quote}
\emph{tā shì yí ge hěn hǎo de yǎnyuán} — He's a very good actor.
\end{quote}

Its characters come later; for now it is a word for the ear.

\begin{cousinweb}
\emph{yǎnyuán} is built from \emph{yǎn} and \emph{yuán}, both inherited from Middle Chinese, \emph{yenX} and \emph{hjwen} (Guangyun, in Baxter's transcription), which glosses the first 'wide' and quotes an older dictionary for the second, 'a number of things'. The senses 'to perform' and 'a member' are not given there, so how the pair came is not settled.
\end{cousinweb}

\begin{grammarlens}[title={What you've built}]
One more word for what I like.
\end{grammarlens}

\subsection*{Your turn}
\begin{itemize}
\raggedright
  \item \emph{Say it:} \emph{yǎnyuán}
  \item \emph{Say it:} \emph{yǎnyuán}, once more
  \item \emph{Say it:} \emph{tā shì yí ge hěn hǎo de yǎnyuán} — He's a very good actor.
  \item \emph{Recall:} say the Mandarin for a girl, then the Mandarin for middle age, then say \emph{yǎnyuán} again
\end{itemize}

\begin{tcolorbox}[breakable,colback=teal!4,colframe=teal!35!black,title={Before you move on}]
What does yǎnyuán mean? (\textquotedblleft{}Actor; performer\textquotedblright{}.) Say it once more.
\end{tcolorbox}
\section[jiémù]{jiémù — programme; TV show}

\label{lesson:ZH-C267-jiemu}



\begin{quote}
Before the new one: say the Mandarin for middle age, then the Mandarin for an actor.
\end{quote}

\subsection*{You'll want to know: jiémù}
\emph{jiémù} — \textquotedblleft{}programme; TV show\textquotedblright{}.

\begin{quote}
\emph{zhège jiémù hěn hǎokàn} — This show is really good.
\end{quote}

Its characters come later; for now it is a word for the ear.

\begin{cousinweb}
\emph{jiémù} is built from \zh{节} and \emph{mù}, both inherited from Middle Chinese, \emph{tset} and \emph{mjuwk} (Guangyun, in Baxter's transcription), which quotes an older dictionary for the first, 'a joint of bamboo', and glosses the second 'the eye'. How the pair came to mean a programme is not settled.
\end{cousinweb}

\begin{grammarlens}[title={What you've built}]
One more word for what I like.
\end{grammarlens}

\subsection*{Your turn}
\begin{itemize}
\raggedright
  \item \emph{Say it:} \emph{jiémù}
  \item \emph{Say it:} \emph{jiémù}, once more
  \item \emph{Say it:} \emph{zhège jiémù hěn hǎokàn} — This show is really good.
  \item \emph{Recall:} say the Mandarin for middle age, then the Mandarin for an actor, then say \emph{jiémù} again
\end{itemize}

\begin{tcolorbox}[breakable,colback=teal!4,colframe=teal!35!black,title={Before you move on}]
What does jiémù mean? (\textquotedblleft{}Programme; TV show\textquotedblright{}.) Say it once more.
\end{tcolorbox}
\section[\texorpdfstring{\zh{明星} (míngxīng)}{明星 (míngxīng)}]{\zh{明星} (míngxīng) — star; celebrity}

\label{lesson:ZH-C267-mingxing}



\begin{quote}
Before the new one: say the Mandarin for an actor, then the Mandarin for a TV show.
\end{quote}

\subsection*{You'll want to know: \zh{明星}}
\textbf{\zh{明星}} — \emph{míngxīng} — \textquotedblleft{}star; celebrity\textquotedblright{}.

\begin{quote}
\textbf{\zh{他是一个大明星}} \emph{tā shì yí ge dà míngxīng} — He's a big star.
\end{quote}

\begin{cousinweb}
\textbf{\zh{明星}} (\emph{míngxīng}) is built from \zh{明} 'bright' and \zh{星} 'a star', both inherited from Middle Chinese, \emph{mjaeng} and \emph{seng} (Guangyun, in Baxter's transcription), which glosses the first 'light' and the second 'a star in the sky'.
\end{cousinweb}

\begin{grammarlens}[title={What you've built}]
One more word for what I like.
\end{grammarlens}

\subsection*{Your turn}
\begin{itemize}
\raggedright
  \item \emph{Say it:} \emph{míngxīng}
  \item \emph{Say it:} \emph{míngxīng}, once more
  \item \emph{Say it:} \emph{tā shì yí ge dà míngxīng} — He's a big star.
  \item \emph{Recall:} say the Mandarin for an actor, then the Mandarin for a TV show, then say \emph{míngxīng} again
\end{itemize}

\begin{tcolorbox}[breakable,colback=teal!4,colframe=teal!35!black,title={Before you move on}]
What does \zh{明星} mean? (\textquotedblleft{}Star; celebrity\textquotedblright{}.) Say it once more.
\end{tcolorbox}
\section[\texorpdfstring{\zh{笑话} (xiàohua)}{笑话 (xiàohua)}]{\zh{笑话} (xiàohua) — joke}

\label{lesson:ZH-C267-xiaohua}



\begin{quote}
Before the new one: say the Mandarin for a TV show, then the Mandarin for a celebrity.
\end{quote}

\subsection*{You'll want to know: \zh{笑话}}
\textbf{\zh{笑话}} — \emph{xiàohua} — \textquotedblleft{}joke\textquotedblright{}.

\begin{quote}
\textbf{\zh{他说了一个笑话}} \emph{tā shuōle yí ge xiàohua} — He told a joke.
\end{quote}

The new character, in pieces: \zh{笑} — ten strokes: two small pieces on top, each a short stroke falling to the left, a short bar and a dot; under them a short stroke falling to the left, then \zh{大}. The next lesson puts it in your hand.

\begin{cousinweb}
\textbf{\zh{笑话}} (\emph{xiàohua}) is built from \zh{笑} 'to laugh' and \zh{话} 'talk', both inherited from Middle Chinese, \emph{sjewH} and \emph{hwaejH} (Guangyun, in Baxter's transcription), which glosses the first 'glad' and the second 'speech'. In the word the second loses its tone.
\end{cousinweb}

\begin{grammarlens}[title={What you've built}]
One more word for what I like.
\end{grammarlens}

\subsection*{Your turn}
\begin{itemize}
\raggedright
  \item \emph{Say it:} \emph{xiàohua}
  \item \emph{Say it:} \emph{xiàohua}, once more
  \item \emph{Say it:} \emph{tā shuōle yí ge xiàohua} — He told a joke.
  \item \emph{Recall:} say the Mandarin for a TV show, then the Mandarin for a celebrity, then say \emph{xiàohua} again
\end{itemize}

\begin{tcolorbox}[breakable,colback=teal!4,colframe=teal!35!black,title={Before you move on}]
What does \zh{笑话} mean? (\textquotedblleft{}Joke\textquotedblright{}.) Say it once more.
\end{tcolorbox}
\section[\texorpdfstring{\zh{笑} (xiào)}{笑 (xiào)}]{\zh{笑} — two small pieces over a short fall and \zh{大}}

\label{lesson:ZH-W267-xiao-laugh}



\begin{quote}
The word first, then the character inside it.

\begin{itemize}
\raggedright
  \item \emph{Recall:} say \emph{xiàohua} — \textbf{R1}, one lesson back, before it decays
\end{itemize}

\begin{quote}
\zh{大} \zh{天}
\end{quote}

\emph{Read it:} each one aloud, with nothing said first — \emph{dà}, \emph{tiān}
\end{quote}

\subsection*{Script — two small pieces over a short fall and \zh{大}}
\hlblockfigure{\detokenize{figures/ZH-W267-xiao-laugh-filmstrip.pdf}}{How \zh{笑} is written, stroke by stroke}

\begin{quote}
\zh{笑}
\end{quote}

\textbf{Two pieces, top before bottom.}

\textbf{The top piece, six strokes: two small pieces side by side,} each three strokes. A short stroke falling to the left, a short bar to the right from it, then a dot under the bar. Then the same three again on the right.

\textbf{The bottom piece, four strokes.} A short stroke falling to the left across the top. Then \zh{大}: a long bar, a stroke falling to the left, and a stroke falling to the right.

6 + 4 = ten.

\subsection*{Your turn}
\begin{itemize}
\raggedright
  \item \emph{Write it:} \zh{笑}, all ten, with the model in front of you
  \item \emph{Cover:} write it again from memory, and count the pieces, not the strokes
\end{itemize}

\begin{tcolorbox}[breakable,colback=teal!4,colframe=teal!35!black,title={Before you move on}]
How many strokes in \zh{笑}? (\textbf{Ten}.) How many small pieces sit on top? (\textbf{Two, three strokes each}.) Which character ends it? (\textbf{\zh{大}}.)

Source: \href{https://raw.githubusercontent.com/chanind/hanzi-writer-data/68d10a4b21150cae5e1ebbd223eed289cf32d90c/data/\%E7\%AC\%91.json}{Hanzi Writer Data \zh{笑}, pinned stroke paths and medians}.
\end{tcolorbox}
\section[\texorpdfstring{\zh{好笑} (hǎoxiào)}{好笑 (hǎoxiào)}]{\zh{好笑} (hǎoxiào) — funny}

\label{lesson:ZH-C267-haoxiao}



\begin{quote}
Before the new one: say the Mandarin for a celebrity, then the Mandarin for a joke.
\end{quote}

\subsection*{You'll want to know: \zh{好笑}}
\textbf{\zh{好笑}} — \emph{hǎoxiào} — \textquotedblleft{}funny\textquotedblright{}.

\begin{quote}
\textbf{\zh{这个笑话很好笑}} \emph{zhège xiàohua hěn hǎoxiào} — This joke is really funny.
\end{quote}

\begin{cousinweb}
\textbf{\zh{好笑}} (\emph{hǎoxiào}) is built from \zh{好} 'good' and \zh{笑} 'to laugh', both inherited from Middle Chinese, \emph{xawX} and \emph{sjewH} (Guangyun, in Baxter's transcription): good to laugh at, the way \zh{好吃} is good to eat.
\end{cousinweb}

\begin{grammarlens}[title={What you've built}]
One more word for what I like.
\end{grammarlens}

\subsection*{Your turn}
\begin{itemize}
\raggedright
  \item \emph{Say it:} \emph{hǎoxiào}
  \item \emph{Say it:} \emph{hǎoxiào}, once more
  \item \emph{Say it:} \emph{zhège xiàohua hěn hǎoxiào} — This joke is really funny.
  \item \emph{Recall:} say the Mandarin for a celebrity, then the Mandarin for a joke, then say \emph{hǎoxiào} again
\end{itemize}

\begin{tcolorbox}[breakable,colback=teal!4,colframe=teal!35!black,title={Before you move on}]
What does \zh{好笑} mean? (\textquotedblleft{}Funny\textquotedblright{}.) Say it once more.
\end{tcolorbox}
